\section{Proofs of the structural propositions}\label{app:proofs}

This appendix collects the full proofs of the three structural
propositions of Section~\ref{sec:prelim}, which rest on classical
network-flow and transportation techniques. The body keeps their
statements, the layered-network construction they build on, and the
facts that the later sections consume, and the proofs appear here in
statement order.

\medskip\noindent
\textbf{Proposition~\ref{prop:oracle} (routing oracle).}
\emph{Fix $(y,z)\in\{0,1\}^{I}\times\{0,1\}^{J}$. The residual routing problem
separates over products, and for each product $l$ its optimal cost equals
the minimum cost $\mu_l$ of an $S$--$T$ flow of value $D_l$ in the layered
network $N_l(y,z)$ of Section~\ref{sec:prelim}, with $\mu_l=+\infty$ when
no such flow exists. The finite $\mu_l$ are attained by integral flows and
are nonnegative integers, and $v(y,z)=\sum_l\mu_l$. Hence $v(y,z)$ is
computable in strongly polynomial time by $|L|$ minimum-cost flow
computations \citep{Orlin1993}.}

\begin{proof}[Proof of Proposition~\ref{prop:oracle}]
We argue through the correspondence between routings and flows. The
routing splits into independent per-product blocks, each block optimum
is identified with the minimum cost of a value-$D_l$ flow in the network
$N_l(y,z)$ constructed in Section~\ref{sec:prelim}, and the claims of the statement all follow
from that identification.

\emph{Separation.} Each of~\eqref{eq:C1}--\eqref{eq:C4} constrains a
single product $l$, and the objective is a sum of per-product terms, so
the feasible set of routings is the product of the per-product feasible
sets and the objective is separable. The minimum is thus the sum over $l$
of the per-product minima, and the design is feasible if and only if every
product admits a routing. In particular, if some product admits no
routing, then its per-product minimum is $+\infty$, the correspondence
established below yields $\mu_l=+\infty$ for that product, and both
$v(y,z)$ and $\sum_l\mu_l$ equal $+\infty$, so the claimed identity holds
under the stated convention. The correspondence itself needs no
finiteness assumption, so we now fix a product $l$ and establish it in
general.

\emph{Routings and flows coincide.} We show the block optimum equals
$\mu_l$. Take a feasible routing $(x,w)$ of block $l$. First lower the
flows into each over-served customer until $\sum_j w_{jkl}=q_{kl}$, then
lower the flows into each depot until $\sum_i x_{ijl}=\sum_k w_{jkl}$.
Which coordinates are lowered is immaterial, and each stopping rule halts
exactly at equality, so \eqref{eq:C1} and~\eqref{eq:C2} hold with
equality afterwards. Feasibility of the routing persists throughout the
lowering. Each lowering decreases only the tracked left-hand side,
\eqref{eq:C1} for $w$ and \eqref{eq:C2} for $x$, which stops exactly at
its target, together with quantities whose decrease merely slackens
\eqref{eq:C2}, \eqref{eq:C3} and~\eqref{eq:C4}. No variable becomes negative, since each lowering drives a
nonnegative total down to a nonnegative target. Each lowering removes
flow from arcs of nonnegative cost, so the cost does not rise.
By~\eqref{eq:C3} at $y_i=0$ and~\eqref{eq:C4} at $z_j=0$, together with
nonnegativity, every $x_{ijl}$ with $y_i=0$ and every $w_{jkl}$ with
$z_j=0$ vanishes, so restricting attention to the open facilities
discards no flow and no cost. Setting $f(S\!\to\! i)=\sum_j x_{ijl}$,
$f(i\!\to\! j_{\mathrm{in}})=x_{ijl}$, the split-arc flow
$=\sum_i x_{ijl}$, $f(j_{\mathrm{out}}\!\to\! k)=w_{jkl}$ and
$f(k\!\to\! T)=q_{kl}$ then conserves flow at every node and respects
every capacity. The result is an $S$--$T$ flow whose value, the total
inflow $\sum_k f(k\!\to\! T)$ at the sink, equals $\sum_k q_{kl}=D_l$,
and its cost is that of the lowered routing, hence at most that of
$(x,w)$.

Conversely, given an $S$--$T$ flow $f$ of value $D_l$, set
$x_{ijl}=f(i\!\to\! j_{\mathrm{in}})$ and
$w_{jkl}=f(j_{\mathrm{out}}\!\to\! k)$ at the open facilities, and set
$x_{ijl}=0$ whenever $y_i=0$ and $w_{jkl}=0$ whenever $z_j=0$. At the
closed facilities \eqref{eq:C3} and~\eqref{eq:C4} read $0\le 0$ and
hold. Value $D_l$ saturates every arc $k\!\to\! T$. Indeed the value
equals $\sum_k f(k\!\to\! T)\le\sum_k q_{kl}=D_l$, so the two sums are
equal, and since each term obeys $f(k\!\to\! T)\le q_{kl}$ every term is
at capacity. Conservation at $k$ then gives $\sum_j w_{jkl}=q_{kl}$,
which is~\eqref{eq:C1}, conservation at the split node gives
$\sum_i x_{ijl}=\sum_k w_{jkl}$, which is~\eqref{eq:C2}, and the split-
and source-arc capacities give~\eqref{eq:C4} and~\eqref{eq:C3} at the
open facilities. The routing so obtained is therefore feasible, and its
cost equals that of the flow because the costs agree arc by arc.
Assembling the two directions, every feasible routing of block $l$
yields a value-$D_l$ flow of no greater cost, so $\mu_l$ is at most the
block optimum, and every value-$D_l$ flow yields a feasible routing of
equal cost, so the block optimum is at most $\mu_l$. The two inequalities
together give the equality, and a value-$D_l$ flow exists if and only if
the block admits a routing, so $\mu_l=+\infty$ exactly when the block is
infeasible.

\emph{Integrality and complexity.} The node--arc incidence matrix of a
flow network is totally unimodular \citep{Schrijver1986}, and all finite
arc capacities $b_{il},p_{jl},q_{kl}$ are integers. Every arc of
$N_l(y,z)$ belongs to one layer, and each layer is an $S$--$T$ cut that
every path crosses exactly once and only forwards, so the arc flows
within a layer sum to the flow value. Since the flows are nonnegative,
every arc in a flow of value $D_l$ therefore carries at most $D_l$, and
the polytope of value-$D_l$ flows is bounded. A finite $\mu_l$ is thus
attained at a vertex of that polytope, since a linear function on a
nonempty bounded polytope attains its minimum at a vertex
\citep{Schrijver1986}. The vertex is integral by the Hoffman--Kruskal
theorem \citep{Schrijver1986}, and the nonnegative integer arc costs
make $\mu_l$ a nonnegative integer. Each $N_l$ has $O(n+|K|)$ nodes and
$O(|I||J|+|J||K|+n+|K|)$ arcs, and a minimum-cost flow of prescribed value is
computed in strongly polynomial time \citep{Orlin1993}. Summing over the
$|L|$ products gives $v(y,z)=\sum_l\mu_l$ within the stated bound.
\end{proof}

\medskip\noindent
\textbf{Proposition~\ref{prop:feas} (feasibility).}
\emph{A design $(y,z)$ is feasible if and only if, for every product $l$,
$\sum_{i\in I}b_{il}\,y_i\ge D_l$ (F1) and
$\sum_{j\in J}p_{jl}\,z_j\ge D_l$ (F2). The test runs in
$O((n+|K|)\,|L|)$ time, of which $O(|K|\,|L|)$ computes the totals
$D_l$.}

\begin{proof}[Proof of Proposition~\ref{prop:feas}]
Necessity will come from aggregating the four constraint groups, and
sufficiency from exhibiting a flow of value $D_l$ in the layered network
$N_l(y,z)$. For necessity fix a feasible routing
and a product $l$.
Summing~\eqref{eq:C1} over $k$ gives $\sum_{j,k}w_{jkl}\ge D_l$, and
summing~\eqref{eq:C4} over $j$ gives
$\sum_{j,k}w_{jkl}\le\sum_j p_{jl}z_j$, so (F2) holds.
Summing~\eqref{eq:C2} over $j$ gives
$\sum_{i,j}x_{ijl}\ge\sum_{j,k}w_{jkl}\ge D_l$, and
summing~\eqref{eq:C3} over $i$ gives
$\sum_{i,j}x_{ijl}\le\sum_i b_{il}y_i$, so (F1) holds.

For sufficiency we route each product $l$ independently and exhibit a flow
of value $D_l$ in $N_l(y,z)$. By the flow-to-routing direction of the
correspondence established for Proposition~\ref{prop:oracle}, any flow of value exactly
$D_l$ in $N_l(y,z)$ reads back as a feasible routing of block $l$, so such
flows for all products assemble into a feasible routing for the design.
The finite-capacity arcs of $N_l(y,z)$ fall into three layers, the source
arcs $S\!\to\! i$ with total capacity $\sum_i b_{il}y_i$, the depot split
arcs with total $\sum_j p_{jl}z_j$, and the demand arcs $k\!\to\! T$ with
total $D_l$. Every $S$--$T$ path crosses each layer exactly once, so each
layer is an $S$--$T$ cut. Since the value of any flow is at most the
capacity of any $S$--$T$ cut \citep{Schrijver1986}, the maximum flow is
at most each of the three totals. We claim it attains the least total,
and since (F1) and (F2) make that least total equal to $D_l$, a flow of
value $D_l$ then exists.

Let $F$ be the least of the three totals.
Because $F$ does not exceed any layer total, one can choose factory loads
$a_i\le b_{il}y_i$, depot loads $e_j\le p_{jl}z_j$ and customer intakes
$r_k\le q_{kl}$ each summing to $F$, for instance by scanning the terms
of each layer in a fixed order and truncating the running sum once it
reaches $F$. Both bipartite stages are complete and uncapacitated, so for
any two nonnegative load vectors with equal totals there is a transfer
between them that meets those loads as its row and column sums. One such
transfer, the northwest-corner procedure of transportation theory,
maintains a current supplier and a current receiver, assigns
them the minimum of the remaining supply and the remaining demand, and
advances past whichever of the two is exhausted. Every step reduces the
residual supply total and the residual demand total by the same amount
and exhausts a supplier or a receiver, so the procedure terminates, and
since the two totals stay equal it ends with every supply and every
demand met exactly, which is the required row and column sums. Route $a$
to $e$ across the first stage and $e$ to $r$ across the second by such
transfers, and put $a_i$ on the source arc $S\!\to\! i$ and $r_k$ on the
demand arc $k\!\to\! T$. Flow is conserved at each factory node because
the first transfer has row sums $a$, and at each customer node because
the second transfer has column sums $r$. Each depot's inflow and outflow
both equal $e_j\le p_{jl}z_j$, so flow is conserved at the split arc and
its capacity is respected. The source and demand arcs carry
$a_i\le b_{il}y_i$ and $r_k\le q_{kl}$, within their capacities, and the
result is an $S$--$T$ flow of value $\sum_k r_k=F$. Hence the maximum
flow equals $F$, and under (F1) and (F2) the least of the three totals is
$D_l$, so a flow of value $D_l$ exists and the design is feasible.
Finally we bound the running time. One pass over the demand matrix
computes every total $D_l$ in $O(|K|\,|L|)$ time, after which the test
evaluates $2|L|$ sums of at most $n$ terms each, for $O((n+|K|)\,|L|)$
in total.
\end{proof}

\medskip\noindent
\textbf{Proposition~\ref{prop:mono} (monotonicity).}
\emph{If $(y,z)\le(y',z')$ componentwise then feasibility of $(y,z)$
implies feasibility of $(y',z')$, and $v(y',z')\le v(y,z)$ always, with
the convention $v=+\infty$ on infeasible designs. In particular the
all-open design attains the least routing value among all feasible
designs.}

\begin{proof}[Proof of Proposition~\ref{prop:mono}]
Since the data
satisfy $b_{il}\ge 0$ and $p_{jl}\ge 0$, the left-hand sides
$\sum_i b_{il}y_i$ and $\sum_j p_{jl}z_j$ of (F1) and (F2) are
nondecreasing in $(y,z)$. If $(y,z)$ is feasible, the necessity direction
of Proposition~\ref{prop:feas} gives (F1) and (F2) at $(y,z)$, hence at
$(y',z')$, and the sufficiency direction then gives feasibility of
$(y',z')$. If $(y,z)$ is infeasible then $v(y,z)=+\infty$ and the
inequality is trivial, so assume $(y,z)$ feasible, in which case
$(y',z')$ is feasible too. Since $(y,z)\le(y',z')$ componentwise, every
node and arc of $N_l(y,z)$ appears in $N_l(y',z')$ with the same capacity
and cost, so $N_l(y,z)$ is a subnetwork of $N_l(y',z')$ and the passage
from one to the other only adds the nodes and arcs of newly opened
facilities. Every $S$--$T$ flow of value $D_l$ in $N_l(y,z)$ is therefore
still one in $N_l(y',z')$, carrying zero on the new arcs, at the same
cost. The feasible set of value-$D_l$ flows only grows and its minimum
cost cannot increase, so $\mu_l(y',z')\le\mu_l(y,z)$ for every $l$, where
$\mu_l(\cdot,\cdot)$ denotes the block optimum of
Proposition~\ref{prop:oracle} viewed as a function of the design. Since
$v=\sum_l\mu_l$ by that proposition, summing over the products gives
$v(y',z')\le v(y,z)$. For the final claim, every design satisfies
$(y,z)\le(\mathbf 1,\mathbf 1)$ componentwise because the variables are
binary. The feasibility part then makes the all-open design feasible
whenever some feasible design exists, and the value part gives
$v(\mathbf 1,\mathbf 1)\le v(y,z)$ for every feasible $(y,z)$, so the
all-open design attains the least routing value among feasible designs,
a claim vacuous when none exists.
\end{proof}
\section{Technical verifications}\label{app:tech}

This appendix collects the proofs whose methods are standard, invariant
tracking for the correctness of the branch-and-bound, linear-programming
duality for the certificates, exchange arguments for the two
preprocessing reductions, dynamic programming for the separable
subclass, and the weight-reduction technique of \citet{FrankTardos1987}
for the compression. The statements remain in the body, each with a
paragraph naming the route of its proof, and the proofs appear here
grouped by result, one subsection each, so that a reference from the
body lands directly on the argument it names.

\subsection{The separable first stage}\label{app:sep}

\noindent
\textbf{Theorem~\ref{thm:separable} (separable first stage).}
\emph{On the group $|K|=|L|=1$, if the first-stage matrix is separable, that is
$c_{ij}=\gamma_i+\delta_j$ for integers $\gamma,\delta$ with the costs
nonnegative, then the problem is solvable in pseudo-polynomial time
$O\bigl(|I||J|+(|I|+|J|)\cdot D\bigr)$.
Separability is decidable in $O(|I||J|)$ time.}

\begin{proof}[Proof of Theorem~\ref{thm:separable}]
The argument reduces the routing cost to a function of its marginals,
and then each facility layer to a knapsack-cover program. With $|K|=|L|=1$ there is
one customer and one product. Write $D=q_{11}$,
$w_j$ for the flow depot $j$ forwards to the customer, and $d_j=d_{j11}$
for its unit second-stage cost. Fix a feasible design $(Y,Z)$. In an
optimal solution we may assume that \eqref{eq:C1} and~\eqref{eq:C2} are
tight, since lowering flow on the nonnegative-cost arcs cannot raise the
cost. The lowering argument is the one in the proof of
Proposition~\ref{prop:oracle} in \ref{app:proofs}. Under
this assumption the depot marginal $e_j=\sum_i x_{ij}=w_j$
is both the inflow and the outflow of depot $j$, the factory marginal
$a_i=\sum_j x_{ij}$ is the load of factory $i$, and $\sum_i a_i=\sum_j
e_j=D$, with $0\le a_i\le b_i$ ($a_i=0$ for $i\notin Y$) and $0\le e_j\le
p_j$ ($e_j=0$ for $j\notin Z$). The total routing cost is
\[
\sum_{ij}c_{ij}x_{ij}+\sum_j d_j w_j
=\sum_{ij}(\gamma_i+\delta_j)x_{ij}+\sum_j d_j e_j
=\sum_i\gamma_i a_i+\sum_j(\delta_j+d_j)\,e_j,
\]
using $\sum_i x_{ij}=e_j$ and $\sum_j x_{ij}=a_i$. The second-stage cost
is thus a per-unit depot cost $d_j$ folded into the depot term. The cost
depends on the transport plan only through the marginals, and on the
complete uncapacitated bipartite stage any nonnegative $a$ and $e$ with
$\sum_i a_i=\sum_j e_j$ are realized by some plan, for instance the
northwest-corner procedure used in the proof of
Proposition~\ref{prop:feas}. A load $a_i>0$ forces $y_i=1$
by~\eqref{eq:C3} and a load $e_j>0$ forces $z_j=1$ by~\eqref{eq:C4},
while an open facility with zero load may be closed at no extra cost
since $f,g\ge 0$, so we may identify the open facilities with the
supports of the loads. Conversely any nonnegative loads within the
capacities and summing to $D$ on each side define a feasible solution
that opens exactly the supports. Adding the fixed charges, the optimum
therefore
splits into two independent problems, one per layer, namely
\[
\min_{a}\Bigl(\sum_{i:a_i>0}f_i+\sum_i\gamma_i a_i\Bigr)
\ \text{ and }\
\min_{e}\Bigl(\sum_{j:e_j>0}g_j+\sum_j(\delta_j+d_j)e_j\Bigr),
\]
each over nonnegative loads bounded by the capacities and summing to $D$.
The two split problems range over real loads while the dynamic program
below ranges over integer ones. For a fixed support the fixed charges
are constant and each problem becomes a linear program over a box
intersected with a single equality, a polytope whose vertices are
integral for integer capacities and integer $D$, so the real and the
integer optima coincide and the restriction to integer loads loses
nothing.
Solve the factory problem by the dynamic program
\[
V(i,s)=\min\Bigl\{\,V(i-1,s),\ \min_{1\le a\le b_i}V(i-1,s-a)+f_i+\gamma_i
a\,\Bigr\},
\]
with $V(0,0)=0$, with $V(0,s)=+\infty$ for $s>0$, with the convention
$V(i-1,s-a)=+\infty$ when $s<a$, and with $s$ ranging over
$0,\dots,D$. The first branch closes factory $i$ and the second opens it
with load $a$. The optimum of the factory problem is $V(|I|,D)$, the
value $+\infty$ signalling that no factory loads reach $D$ within the
capacities, in which case the instance is infeasible. Substituting
$t=s-a$ linearizes the inner minimum, since
\[
\min_{1\le a\le b_i}V(i-1,s-a)+f_i+\gamma_i a
= f_i+\gamma_i s+\min_{\max(0,s-b_i)\le t\le s-1}
\bigl(V(i-1,t)-\gamma_i t\bigr).
\]
As $s$ runs from $0$ to $D$ the window of admissible $t$ slides
monotonically to the right, so a monotone double-ended queue over the
keys $V(i-1,t)-\gamma_i t$ returns each window minimum in $O(1)$
amortized time, every index entering and leaving the queue once per
row. The device works whatever the sign of $\gamma_i$, and the
infinite entries are skipped since they can never attain a finite
minimum. Each row therefore
costs $O(D)$ and the factory table $O(|I|\cdot D)$ time. The depot
problem is the
same recurrence with $g_j$, $\delta_j+d_j$ and $p_j$, in $O(|J|\cdot
D)$ time, its value $+\infty$ likewise signalling
infeasibility. Their sum is the instance optimum, within the
stated bound. The matrix is separable exactly when
$c_{ij}+c_{11}=c_{i1}+c_{1j}$ for all $i,j$, the condition being necessary by
substituting $c_{ij}=\gamma_i+\delta_j$, and sufficient because then
$\gamma_i=c_{i1}-c_{11}$ and $\delta_j=c_{1j}$ satisfy
$\gamma_i+\delta_j=c_{i1}+c_{1j}-c_{11}=c_{ij}$, and it is checked in
$O(|I||J|)$ time. The $O(|I||J|)$ term of the stated running time reads
the matrix, checks separability and recovers $\gamma$ and $\delta$.
\end{proof}

\subsection{Correctness of the branch-and-bound}\label{app:bb}

\noindent
\textbf{Proposition~\ref{prop:bbcorrect} (correctness).}
\emph{Algorithm~\ref{alg:xpbb} decides $\mathrm{MP\text{-}TSCFLP}(B,k)$, and when
run with $B=+\infty$ it returns $\mathrm{best}=\OPT_k$.}

\begin{proof}[Proof of Proposition~\ref{prop:bbcorrect}]
The proof runs in three steps, justifying the early return of line 3,
bounding $\mathrm{best}$ from below by soundness, and bounding it from
above by tracking one distinguished optimum through the search.

\emph{The early return.} Line 3 is the root instance of P1. If some
feasible design $O$ opens at most $k$ facilities, then for every product
$l$ the factories of $O$ satisfy (F1), so Lemma~\ref{lem:cover} gives
$k^\ast_I(l)\le|O\cap I|$, and likewise $k^\ast_J(l)\le|O\cap J|$, whence
$\max_l k^\ast_I(l)+\max_l k^\ast_J(l)\le|O|\le k$ and the return does
not fire. When it fires, therefore, no feasible design opens at most $k$
facilities, so $\OPT_k=+\infty$ and the answer \textsc{no} is correct.

\emph{Soundness.} $\mathrm{best}$ is assigned only at a leaf, to
$LB(O,C)=\mathrm{cost}(O)$. Such a leaf survived P2, so
$v^{\uparrow}=v(O\cup\emptyset)=v(O)<+\infty$ and $O$ is feasible, and it
survived the cardinality test of line 6, so $r=k-|O|\ge 0$ and
$|O|\le k$. Hence $\mathrm{best}\ge\OPT_k$ throughout, where
$\OPT_k=+\infty$ if no feasible design opens at most $k$ facilities.

\emph{Completeness under $\OPT_k\le B$.} Suppose $\OPT_k\le B$ and fix
$O^\ast$, a feasible design with
$|O^\ast|\le k$ and $\mathrm{cost}(O^\ast)=\OPT_k$, of least cardinality
among such. Call a stack node $(O,C)$ \emph{live} if $O\subseteq O^\ast$
and $O^\ast\cap C=\emptyset$, in which case $O^\ast\subseteq O\cup F$,
since $O^\ast$ avoids $C$ and every facility lies in $O\cup C\cup F$. A
step of the search is one iteration of the while loop, and we show by
induction over steps that either $\mathrm{best}=\OPT_k$ or the stack
holds a live node. Before the first step the stack holds the root, which
is live. Once $\mathrm{best}=\OPT_k$ holds it holds forever, since
$\mathrm{best}$ never increases and soundness bounds it below by
$\OPT_k$. When a step pops a node that is not live while a live node is
on the stack, the live node stays there untouched and the invariant
survives, so assume the popped node $(O,C)$ is live and
$\mathrm{best}>\OPT_k$.

We first check that no test removes the live node. It has
$|O|\le|O^\ast|\le k$, so $r\ge 0$ and line 6 keeps it. Beyond $O$ the
design $O^\ast$ opens only free facilities, and for every product $l$
the added factories $(O^\ast\setminus O)\cap I$ supply the residual
demand of (F1), so Lemma~\ref{lem:cover} gives
$s_I(l;O,C)\le|(O^\ast\setminus O)\cap I|$ and likewise
$s_J(l;O,C)\le|(O^\ast\setminus O)\cap J|$. Hence
\[
\max_l s_I(l;O,C)+\max_l s_J(l;O,C)\ \le\ |O^\ast\setminus O|
\ =\ |O^\ast|-|O|\ \le\ k-|O|\ =\ r,
\]
and P1 spares the node. Monotonicity along $O^\ast\subseteq O\cup F$
gives $v^{\uparrow}=v(O\cup F)\le v(O^\ast)<+\infty$
(Proposition~\ref{prop:mono}), so the infeasibility branch of P2 does
not fire, and $O^\ast$ is a completion of $(O,C)$, so
$LB(O,C)\le\mathrm{cost}(O^\ast)=\OPT_k<\mathrm{best}$, which with
$\OPT_k\le B$ gives $LB(O,C)<\min(\mathrm{best},B+1)$ and the bound
branch of P2 spares it as well.

If the live node is internal it is expanded on a free facility $e$, and
the child that records the status of $e$ in $O^\ast$, namely
$(O\cup\{e\},C)$ when $e\in O^\ast$ and $(O,C\cup\{e\})$ otherwise, is
live, so the invariant survives. If the live node is a leaf then
$F=\emptyset$, so $O^\ast\subseteq O\cup F=O$, and $O\subseteq O^\ast$
forces $O=O^\ast$. The oracle solution returned at the leaf is an
optimal flow of $O^\ast$ and leaves no open facility
unused, since closing an unused facility would give a feasible design
opening at most $k$ facilities at cost at most $\OPT_k$, hence an
optimum of smaller cardinality than $O^\ast$, against its choice. P3
therefore does not fire, and the leaf sets $\mathrm{best}$ to
$LB(O^\ast,C)=\mathrm{cost}(O^\ast)=\OPT_k$, after which the invariant
persists in its first form. Each child carries one fewer free facility
than its parent, so the search tree has depth at most $n$ and finitely
many nodes, and the loop terminates. At termination the stack is empty,
so the invariant leaves only $\mathrm{best}=\OPT_k$.

\emph{Conclusion.} If $\OPT_k\le B$ then completeness gives
$\mathrm{best}=\OPT_k\le B$, a finite value since the budget $B$ of the
decision problem is a nonnegative integer, so both tests of line 14 pass
and the algorithm returns \textsc{yes}. If $\OPT_k>B$ then
$\mathrm{best}\ge\OPT_k>B$ by soundness, one of the two tests fails
whether $\mathrm{best}$ is finite or not, and the algorithm returns
\textsc{no}, and the answer \textsc{no} of line 3, when that return
fires instead, is correct by the early-return argument. Either
way the algorithm decides $\mathrm{MP\text{-}TSCFLP}(B,k)$. Running with
$B=+\infty$ makes $\OPT_k\le B$ hold whenever $\OPT_k$ is finite, so
completeness gives $\mathrm{best}=\OPT_k$ in the finite case. When
$\OPT_k=+\infty$ soundness keeps $\mathrm{best}$ unassigned, so it
retains its initial value $+\infty=\OPT_k$, the finiteness test of
line 14 fails and the label is \textsc{no}, and when the early return
fires instead line 3 reports the same infinite optimum. On either
branch of line 14 the reported $\mathrm{best}$ equals $\OPT_k$, so the
run with $B=+\infty$ reads the optimum off every exit.
\end{proof}

\subsection{Routing certificates and Benders cuts}\label{app:duals}

\noindent
\textbf{Proposition~\ref{prop:duals} (routing certificate and Benders cuts).}
\emph{Fix a product $l$ and let $\Delta_l$ be the dual polyhedron of its
transportation block, which is independent of $(y,z)$. If the block is
feasible then some $u^\ast\in\Delta_l$ attains the dual value
$\mu_l(y,z)$, and the pair of an optimal flow and $u^\ast$ certifies
$\mu_l(y,z)$ in time $O(|I||J|+|J||K|)$. Further, writing $\theta_l$ for
a variable through which a master problem represents the block value, for
every $u=(\alpha,\beta,\sigma,\pi)\in\Delta_l$ the inequality
\[
\theta_l\ \ge\ \sum_k q_{kl}\,\alpha_k
      -\sum_i b_{il}\,\sigma_i\,y_i-\sum_j p_{jl}\,\pi_j\,z_j
\]
is valid for every design whose block $l$ is feasible, and it holds with
equality at the design that generated $u^\ast$. Two explicit recession
rays of $\Delta_l$ certify infeasibility, and their objectives are
positive exactly when (F1), respectively (F2), of
Proposition~\ref{prop:feas} fails.}

\begin{proof}[Proof of Proposition~\ref{prop:duals}]
The block dual, once written explicitly, yields the certificate, the
cuts and the rays. Write the block primal as the minimization
of
$\sum_{ij}c_{ijl}x_{ijl}+\sum_{jk}d_{jkl}w_{jkl}$ over nonnegative flows
subject to the demand rows~\eqref{eq:C1}, the balance
rows~\eqref{eq:C2}, oriented as $\sum_i x_{ijl}-\sum_k w_{jkl}\ge 0$,
the factory rows~\eqref{eq:C3} with right-hand side $b_{il}y_i$, and the
depot rows~\eqref{eq:C4} with right-hand side $p_{jl}z_j$. The trimming
step in the proof of Proposition~\ref{prop:oracle}, given in
\ref{app:proofs}, lowers flows at no
cost increase until~\eqref{eq:C2} is tight, so replacing~\eqref{eq:C2}
by the equality $\sum_i x_{ijl}=\sum_k w_{jkl}$ changes neither the
feasibility nor the optimal value of the block, and we dualize it in
that equality form, which is what licenses a free multiplier on the
balance rows. Assign dual
multipliers $\alpha_k\ge 0$ to the demand rows, a free $\beta_j$ to the
balance equalities, and $\sigma_i\ge 0$ and
$\pi_j\ge 0$ to the two capacity groups. The dual constraints, one per primal
variable, read $\beta_j-\sigma_i\le c_{ijl}$ for $x_{ijl}$ and
$\alpha_k-\beta_j-\pi_j\le d_{jkl}$ for $w_{jkl}$, and none of them
mentions the design, so the dual feasible set $\Delta_l$ is indeed
independent of $(y,z)$. The dual objective is
$\sum_k q_{kl}\alpha_k-\sum_i b_{il}y_i\sigma_i-\sum_j p_{jl}z_j\pi_j$,
so weak duality gives, for every $u\in\Delta_l$, the displayed lower
bound on the block value $\theta_l$, valid across all designs whose block
is feasible precisely because $\Delta_l$ carries no design. When the block
is feasible its primal is bounded below by $0$ through the nonnegative
costs, so the finite optimum makes linear-programming strong duality
supply a $u^\ast\in\Delta_l$ attaining $\mu_l(y,z)$, giving equality
there. Checking a certificate pair is verifying primal and dual
feasibility and equal objectives, which certifies optimality by weak
duality, in one pass over the $O(|I||J|+|J||K|)$ coefficients. When
the block is infeasible the primal is empty, and $\Delta_l$ contains the
point $u=0$ by nonnegativity of the costs $c$ and $d$, so the dual is
feasible and therefore unbounded
along a recession ray of $\Delta_l$, whose cone is $\beta_j-\sigma_i\le
0$, $\alpha_k-\beta_j-\pi_j\le 0$ with $\alpha,\sigma,\pi\ge 0$. Two
rays realize the feasibility cuts. The ray $\alpha\equiv 1$, $\beta\equiv
1$, $\sigma\equiv 1$, $\pi\equiv 0$ lies in the cone, since its constraint
values are $\beta_j-\sigma_i=1-1=0\le 0$ and
$\alpha_k-\beta_j-\pi_j=1-1-0=0\le 0$, and it has objective
$\sum_k q_{kl}-\sum_i b_{il}y_i=D_l-\sum_i b_{il}y_i$, positive exactly
when (F1) fails. The ray $\alpha\equiv 1$, $\beta\equiv 0$, $\sigma\equiv
0$, $\pi\equiv 1$ lies in the cone through $\beta_j-\sigma_i=0\le 0$
and $\alpha_k-\beta_j-\pi_j=1-0-1=0\le 0$, and has objective
$D_l-\sum_j p_{jl}z_j$, positive exactly
when (F2) fails. For the concluding equivalence note that block $l$ is
feasible exactly when (F1) and (F2) hold for $l$, since both the
aggregation argument for necessity and the flow construction for
sufficiency in the proof of Proposition~\ref{prop:feas} treat each
product on its own. A design's block is therefore infeasible if and only
if (F1) or (F2) fails for $l$, which happens if and only if one of
these two rays has positive objective, recovering the aggregate test of
Proposition~\ref{prop:feas}.
\end{proof}

\subsection{Customer aggregation}\label{app:agg}

\noindent
\textbf{Proposition~\ref{prop:aggregation} (customer aggregation).}
\emph{If customers $k\ne k'$ have identical cost columns, that is
$d_{jkl}=d_{jk'l}$ for all $j,l$, then merging $k'$ into $k$ and adding
their demands preserves feasibility, routing value and cost of every
design, hence the optimal value, the set of optimal designs and the
answers of both decision versions. Feasible complete solutions
map in both directions under the explicit merge and split maps of the
proof, which never raise the cost and preserve the optimal routing
value of every design, and consequently one
may assume $|K|\le\tau$, where $\tau$
is the number of distinct second-stage cost columns, computable in
polynomial time.}

\begin{proof}[Proof of Proposition~\ref{prop:aggregation}]
Fix a design and a product $l$. The first stage and the fixed
charges are
untouched, so only the second-stage cost of serving $k$ and $k'$ against
the merged customer $\hat k$ with demand $q_{kl}+q_{k'l}$ is at issue,
and the two customers share the cost column $d_{j\bullet l}$. Given a
routing serving both, the merge map sets
$w_{j\hat kl}=w_{jkl}+w_{jk'l}$. This serves
the merged demand, leaves every depot total unchanged, so~\eqref{eq:C2}
and~\eqref{eq:C4} still hold, and costs
$\sum_j d_{j\bullet l}(w_{jkl}+w_{jk'l})$, exactly the original cost.
Conversely, given a routing serving $\hat k$, trim any surplus so that
$\sum_j w_{j\hat kl}=q_{kl}+q_{k'l}$, which only slackens~\eqref{eq:C2}
and~\eqref{eq:C4} and, by nonnegativity of $d$, does not raise the cost,
then apply the split map, a
northwest-corner expansion of each $w_{j\hat kl}$ into $w_{jkl}+w_{jk'l}$ with
$\sum_j w_{jkl}=q_{kl}$ and $\sum_j w_{jk'l}=q_{k'l}$, which is possible
because only the totals are constrained. The cost is again unchanged.
The two instances therefore have the same feasible designs with the same
routing values and costs, hence the same optimal value, the same
optimal designs and the same decision answers. The merge map preserves
the cost exactly, and the split map, run after the trim, never raises
it, so the two maps carry feasible complete solutions across in both
directions and preserve the optimal routing value of every design,
although the trim may strictly lower the cost of a routing that
carries surplus, so the correspondence is not a cost-preserving
bijection on all complete solutions. Iterating over
equal columns leaves at most one customer per
distinct column, so $|K|\le\tau$.
\end{proof}

\subsection{Capacity capping}\label{app:cap}

\noindent
\textbf{Observation~\ref{obs:capping} (capacity capping).}
\emph{Replacing every capacity $b_{il}$ by $\min(b_{il},D_l)$ and every
$p_{jl}$ by $\min(p_{jl},D_l)$ preserves feasibility, routing value and
cost of every design, hence the optimum and both decision answers.}

\begin{proof}[Proof of Observation~\ref{obs:capping}]
Fix a design of
the original instance. Any routing of it can be trimmed, as in
the proof of Proposition~\ref{prop:oracle} in \ref{app:proofs}, until each product's total
flow across each stage equals $D_l$, at no increase in cost. A
trimmed routing sends at most $D_l$ units of product $l$ through any
single factory or depot, so it respects the capped capacities and is a
routing of the capped instance at the same cost. Conversely every
routing of the capped instance is one of the original, because capping
only lowers capacities, so the two instances give every design the same
routing value. Feasibility is also unchanged, since
$\sum_i\min(b_{il},D_l)\,y_i\ge D_l$ holds if and only if
$\sum_i b_{il}\,y_i\ge D_l$, either because some open factory alone
reaches $D_l$ or because no cap binds, and the depot side is identical.
The two instances therefore share feasible designs, routing values and
costs, hence the optimum and both decision answers.
\end{proof}

\subsection{Kernel and compression under one-sided bounds}\label{app:ft}

\noindent
\textbf{Proposition~\ref{prop:ftcompress} (kernel and compression
under one-sided bounds).}
\emph{Fix $c^\star$ and let $\kappa=|I|+|J|+|L|+\tau$. The subclass of
$\mathrm{MP\text{-}TSCFLP}(B)$ whose opening and transport costs
$f,g,c,d$ and whose total demands $D_1,\dots,D_{|L|}$ are integers at
most $\kappa^{c^\star}$, with capacities and budget unrestricted,
admits a polynomial kernel in $\kappa$, an equivalent instance of
$\mathrm{MP\text{-}TSCFLP}(B)$ itself of $\mathrm{poly}(\kappa)$
bits. The subclass with $\max(f,g,c,d,B)\le\kappa^{c^\star}$ and
capacities and demands unrestricted admits a
polynomial compression in $\kappa$ into a fixed sign-condition
language, which is not itself a location problem. The first subclass is nontrivial.
The strongly hard slice with $|K|=|L|=1$ of Theorem~\ref{thm:numeric}
has costs in $\{0,1\}$ and a single customer whose demand, the only
product total, is $|I|\,(|J|+1)\le\kappa^{2}$, so the slice lies in the
subclass for every $c^\star\ge 2$.}

\begin{proof}[Proof of Proposition~\ref{prop:ftcompress}]
Consider the first subclass, in which the opening and transport costs
and the totals $D_l$ are integers at most $\kappa^{c^\star}$ while the
capacities and the budget carry no bound.
Proposition~\ref{prop:aggregation} leaves $|K|\le\tau$ customers, each
merged demand entry is a sum of original entries of the same product
and hence at most $D_l\le\kappa^{c^\star}$, and the totals $D_l$
themselves survive the merge unchanged. The capping of
Observation~\ref{obs:capping} then bounds every capacity by
$\max_l D_l\le\kappa^{c^\star}$. After these two reductions at most
$\kappa$ facilities, $\tau\le\kappa$ customers and $|L|\le\kappa$
products remain and every datum except the budget is an integer at
most $\kappa^{c^\star}$, so the instance without its budget occupies
$\mathrm{poly}(\kappa)$ bits and only $B$ stands in the way of a
kernel.

The budget is capped against the all-open design. If that design is
infeasible then no design is feasible, since every design lies
componentwise below it and feasibility is monotone by
Proposition~\ref{prop:mono}, and the kernel outputs a fixed
no-instance, a single facility pair of zero capacities facing one unit
of demand. Otherwise one oracle call prices the all-open design in
polynomial time (Proposition~\ref{prop:oracle}) at
$U=f(I)+g(J)+v(\mathbf 1,\mathbf 1)$. The fixed charges total at most
$\kappa\cdot\kappa^{c^\star}$, at most
$\sum_l D_l\le\kappa^{c^\star+1}$ units of flow move in all, and each
unit traverses one first-stage and one second-stage arc of costs at
most $\kappa^{c^\star}$ each, so
$U\le\kappa^{c^\star+1}+2\kappa^{2c^\star+1}\le
3\kappa^{2c^\star+1}$. If $B\ge U$ the all-open design already
witnesses acceptance and the kernel outputs a fixed yes-instance, one
facility pair of zero costs and zero demand with budget zero. Otherwise
$B<U\le 3\kappa^{2c^\star+1}$, the budget fits in $O(\log\kappa)$
bits for fixed $c^\star$, and the reduced instance, an ordinary
instance of $\mathrm{MP\text{-}TSCFLP}(B)$ equivalent to the original
one, occupies $\mathrm{poly}(\kappa)$ bits in total, which is the
kernel claimed.

In the second subclass the costs and the budget are at most
$\kappa^{c^\star}$ and the large numbers are the capacities and
demands, so no polynomially bounded cost of a witness design is
available to cap them and the argument changes tool. The compression
treats the extended vector
$\Psi=(b,p,q,1)$ at once, whose final coordinate absorbs the additive
constants below, and customer aggregation brings the dimension of
$\Psi$ down to $\mathrm{poly}(\kappa)$. For any fixed design $S$ the
routing value is the sum over products $l$ of a minimum-cost flow in
the network $N_l(S)$ of Proposition~\ref{prop:oracle}, and each
per-product minimum is attained at a basic (vertex) flow. The
node--arc incidence matrix of each $N_l(S)$ is totally unimodular
\citep{Schrijver1986}, so every subdeterminant lies in $\{-1,0,1\}$, a
nonsingular basis matrix $M$ has determinant $\pm 1$, and by Cramer's
rule each entry of $M^{-1}$, a cofactor divided by that determinant,
lies in $\{-1,0,1\}$ as well. Every basic-flow value is therefore a
$\{-1,0,1\}$-combination of right-hand sides drawn from $\Psi$, so the
nonnegativity of a basic flow and the comparison of its routing cost
against the budget $B$ minus the fixed charge of the design are sign
conditions on integer linear forms in $\Psi$. The coefficients of these
forms on the capacity--demand entries are signed sums of the small cost
entries, of magnitude $\mathrm{poly}(\kappa)$, the coefficient on the
final entry is an integer of magnitude at most $B$ plus the total fixed
charge, again $\mathrm{poly}(\kappa)$, and the tests (F1) and (F2)
contribute further forms with coefficients in $\{-1,0,1\}$. Acceptance
is thus a Boolean combination of sign conditions on integer linear
forms in $\Psi$ of dimension $\nu=\mathrm{poly}(\kappa)$ with
coefficients bounded by some
$\Delta_{\max}=\mathrm{poly}(\kappa)$. The Frank--Tardos theorem
\citep{FrankTardos1987,EtscheidKratschMnichRoeglin2017} produces, for
any prescribed bound $N_0$, in time
polynomial in the encoding length of $\Psi$ and in $\log N_0$, an integer vector $\Psi'$ of
$\mathrm{poly}(\nu,\log N_0)$ bits such that
$\operatorname{sign}(a\cdot\Psi')=\operatorname{sign}(a\cdot\Psi)$ for
every integer vector $a$ whose $\ell_1$-norm is strictly below $N_0$.
We invoke it with $N_0=\nu\,\Delta_{\max}+1$. Every acceptance form
has $\lVert a\rVert_{\infty}\le\Delta_{\max}$ and hence
$\lVert a\rVert_{1}\le\nu\,\Delta_{\max}=N_0-1$, so every form is
strictly covered, and $\log N_0=O(\log\kappa)$ for fixed $c^\star$,
so $\Psi'$ has $\mathrm{poly}(\kappa)$ bits and replacing $\Psi$ by
$\Psi'$ preserves the answer. The output pairs the unchanged
small data with the compressed vector, the Boolean combination of sign
conditions still decides acceptance, and the map is a polynomial
compression onto the fixed language of these sign conditions rather
than a well-formed instance of $\mathrm{MP\text{-}TSCFLP}(B)$, because
entries of $\Psi'$ may be negative. For each fixed $c^\star$ that
target language is a fixed problem, and the compression framework
places no requirement that it lie in \NP. With
both sides free the acceptance condition compares products of a large cost
with a large capacity, the forms become bilinear, and the tool no longer
applies.
\end{proof}
